% Model & Report, Part I.  Built by `just` (see the Makefile).  Numbers come from build/report/numbers.tex,
% which scripts/make_report.py writes from results/results.json.  Do not type results into this file.
\documentclass[11pt]{article}
\usepackage[margin=1in]{geometry}
\usepackage{graphicx}
\usepackage{booktabs}
\graphicspath{{results/figures/}}
\input{build/report/numbers.tex}

\newcommand{\reportauthor}{Janice Yeung}   % change here if needed
\setlength{\parindent}{0pt}
\setlength{\parskip}{0.6em}

\title{Time to failure of IDC sensors from before/after pictures\\[0.3em]
       \large EMA 5937 --- Model \& Report, Part I (pipeline milestone)}
\author{\reportauthor}
\date{}

\begin{document}
\maketitle

\section{What this report is}

This milestone checks that the whole chain works: database $\rightarrow$ pictures $\rightarrow$ model
$\rightarrow$ held-out scores $\rightarrow$ figures $\rightarrow$ this report, all from one command (\texttt{just}).
The model is a deliberately simple baseline; making it good is Part II.  The main result is therefore
the \emph{structure} of the evaluation: a held-out score shown next to a negative control (a dummy
that cannot learn) and a positive control (a ``model'' that is given the answer).

\section{Data and split}

One sample is one sensor (U1--U4) on one board with a measured time to failure.  Its input is two
grey-scale pictures of the sensor, pristine and exposed, each shrunk to \Pixels$\times$\Pixels\ pixels
(from the tensors built at scale \Scale), plus which sensor it is.  The quantity to predict is
$\log_{10}$ of the time to failure in milliseconds.

All sensors of a board go to the same side of the split, so the model is never tested on a board it
has seen.  The training file has \NTrainFile\ sensors from \NBoardsTrain\ boards.  Of these,
cross-validation fold \ValFold\ (\NCheck\ sensors, whole boards) is the \emph{check set} used to watch training, and
the other \NLearned\ sensors are learned from.  The \emph{held-out test set} is a separate file of
\NTest\ sensors from \NBoardsTest\ boards (\PerSensor); the number of boards shared with training
is \NBoardsShared.  The test file is opened only by the scoring step, after training has finished.

\section{Model and training}

The model is a shallow neural network written in numpy: \Pixels$\times$\Pixels$\times$2 picture
values and a one-hot sensor label go in, one hidden layer of \Hidden\ units follows, and one number
(the guessed $\log_{10}$ time) comes out (\NWeights\ weights in total).  It is trained with Adam (learning rate
\LearningRate, weight decay \WeightDecay, batches of \Batch) for \Epochs\ epochs (\NSteps\ steps), and the weights
from the epoch with the lowest check-set loss are kept (epoch \BestEpoch).

Figure~\ref{fig:train} shows the loss on the training samples and on the check set.  A value of 1 means
no better than always guessing the average.  A training loss that keeps falling while the check-set loss
stalls or rises means the model is starting to memorise rather than learn, which is why the best epoch is chosen on
the check set.  Here the check-set loss at the best epoch is \BestValLoss, against \DummyValLoss\ for the dummy,
and the training loss at the last epoch is \FinalTrainLoss.

\begin{figure}[htbp]
\centering
\includegraphics{training_curve.pdf}
\caption{Training curve.  Dashed orange line: the dummy on the same check set.}
\label{fig:train}
\end{figure}

\section{Controls}

\textbf{Negative control (dummy).}  Ignores the pictures and always guesses the average
$\log_{10}$ time of the training set.  It cannot learn, so its AUC must be exactly 0.5 and its $R^2$ at or below 0.
A model that does not beat this has learned nothing.

\textbf{Positive control.}  Cheats on purpose.  The last time stamp of a sensor's current-versus-time curve
\emph{is} the time to failure (a label leak recorded in the ETL notes), so a straight line through it,
fitted on the training samples, should predict almost perfectly.  It is not a usable model, because the
current curve is only known once the sensor has failed.  Its job is to show that the loading, ordering and
scoring code can detect a model that has the signal.  Only \NWithCurrent\ of the \NTest\ test sensors
have a current file; the rest get the training average, which limits its overall score.  Table~\ref{tab:subset}
therefore also scores all three on just the sensors that have a current file.

\section{Held-out results}

Scores on the held-out test set (\NTest\ sensors).  $R^2$ is 1 for perfect guesses and 0 for always guessing the
test-set average.  ``Typical miss'' is $10^{\mathrm{MAE}}$: a value of 2$\times$ means the guess is usually within a factor
of two of the true time.  AUC asks whether the guess separates short-lived sensors (below the training median of
\ThresholdSeconds\ s; \NShort\ of the \NTest\ test sensors) from long-lived ones; 0.5 is a coin flip.  The 95\% intervals
re-draw whole boards \NBoot\ times, because sensors on one board are not independent.

\begin{table}[htbp]
\centering
\small
\metricstable
\caption{Held-out test scores (all \NTest\ sensors).}
\label{tab:main}
\end{table}

\AucVerdict

\begin{table}[htbp]
\centering
\small
\subsettable
\caption{The same scores on only the \NWithCurrent\ test sensors that have a current file.}
\label{tab:subset}
\end{table}

\begin{figure}[htbp]
\centering
\includegraphics{roc.pdf}
\caption{ROC curves on the held-out test set.  The dummy lies on the diagonal.}
\label{fig:roc}
\end{figure}

\begin{figure}[htbp]
\centering
\includegraphics{parity.pdf}
\caption{Guessed versus true $\log_{10}$ time to failure on the held-out test set.  The grey line is a perfect guess.
Flat rows of dots are samples where the guess is the training average (always for the dummy; for the positive
control, sensors without a current file).}
\label{fig:parity}
\end{figure}

\section{What this does and does not show}

\begin{itemize}
\item The chain runs end to end, and changing the data changes the figures and tables: the Makefile
      rebuilds the database, the tensors, the model, the scores, the figures and this report in that order.
\item The controls are checked by \texttt{just check}: the dummy must score AUC 0.5 and $R^2$ at or below 0, and the positive
      control must reproduce the truth on the sensors that have a current file.  In this run the dummy scores AUC \AucDummy\ and
      $R^2$ \RsqDummy, and the positive control scores $R^2$ \RsqSubPos\ and AUC \AucSubPos\ on those sensors.
\item The test set is small (\NBoardsTest\ boards), so the intervals are wide and the model's score
      should be read as ``consistent with'' rather than ``shown to be''.  Also, \NShort\ of the \NTest\ test sensors lie below the
      training median; a share far from one half means the test boards differ from the training boards, which
      lowers $R^2$ even when the ranking (AUC) is reasonable.
\item The pictures may carry information about the test conditions or the scan batch (for example the applied
      voltage, which is not an input) rather than about degradation.  Part II should check what the model is
      actually using.
\end{itemize}

\section{Next steps (Part II)}

\emph{To be written for Part II:} the model choice (for example DINOv3 or U-Net), what changes in the data, how the
score is optimised (on the check set or by cross-validation, never on the test set), and the discussion of results.

\end{document}
